\documentclass[10pt,a4paper]{amsart}
\usepackage[margin=19mm]{geometry}
\usepackage{amsmath,amssymb,mathtools}
\usepackage[scr=rsfs,cal=euler]{mathalfa}
\usepackage{xfrac}
\usepackage[unicode,hidelinks,bookmarks=false]{hyperref}
\newcommand{\ie}{i.e.\ }
\newcommand{\cf}{cf.\ }
\newcommand{\resp}{respectively\ }
\newcommand{\Subsection}{\subsection}
\providecommand{\nobreakdash}{\nobreak\hbox{-}}
\setlength{\emergencystretch}{2em}
\multlinegap0pt
\usepackage{bm}
\usepackage{bbm}
\usepackage{dsfont}
\usepackage{xspace}
\usepackage{cancel}
\usepackage{mathtools}
\usepackage{amsthm}
\makeatletter
\@ifundefined{Theorem}{\newtheorem{Theorem}{Theorem}}{}
\@ifundefined{Lemma}{\newtheorem{Lemma}[Theorem]{Lemma}}{}
\makeatother
\newcommand{\R}{\mathcal{R}}
\title[Motives of Nullcones of Quiver Representations]{Motives of Nullcones of Quiver Representations}
\author{Lydia G\"osmann, Markus Reineke}
\date{Source: arXiv:2502.02353v3 (CC BY 4.0)}
\begin{document}
\maketitle

\noindent\emph{Source and licence.} Motives of Nullcones of Quiver Representations. Version arXiv:2502.02353v3 (CC BY 4.0), distributed under the Creative Commons Attribution 4.0 International licence, as recorded in the arXiv licence metadata for this exact version and shown on the abstract page https://arxiv.org/abs/2502.02353v3. Full rights notice: licenses/arxiv-2502.02353-CC-BY-4.0.txt.\par\smallskip

\noindent\textit{Editorial scope.} Counted unit: the Theorem whose source label is \texttt{theo balanced semi-stable empty} in the article (source0.tex lines 571--573, source label \texttt{theo balanced semi-stable empty}), delivered together with the article's own complete original proof of it (source0.tex lines 575--618, one \texttt{proof} environment of 44 lines). No mathematics is written, repaired or re-derived: statement and proof are the article's own text line for line. Required context retained uncounted: lem inequ for left-complete (lines 563--567). Every definition, display and label of those lines is retained unchanged. Omitted from this package and not redistributed: 1--562, 568--570, 574--574, 619--981. Nothing inside a retained excerpt is omitted or altered: every retained body is delivered exactly as the article's lines are written, so no omission receipt of any kind accompanies this package. Citations: the retained excerpts carry no citation command at all, so no bibliography entry of the article is transcribed and no reference is redistributed. Reference adaptation: none was needed and none was made; every reference inside the delivered unit resolves inside it, so no unresolved reference or citation is produced. Printed slips kept literal: no printed word, symbol, hypothesis or number of the counted statement or of its proof was altered. Layout and class: the article's own class is \texttt{sigma} with packages including ifpdf, amscd, fancyhdr, stmaryrd, bbm, mathtools, tikz-cd, xy, extarrows; the wrapper is the lane's pinned \texttt{amsart} editorial wrapper at 10pt on A4, which restates the theorem-like environments the retained text needs and adds the article's own macro definitions that the retained text needs (\texttt{\textbackslash R}), with extra wrapper packages (bm, bbm, dsfont, xspace, cancel, mathtools). The delivered numbering is the unit's own. Assets: no retained span contains a figure environment, a graphics-inclusion command, a PDF-page inclusion, a table or a TikZ picture; no image file is copied into this package, the package declares no asset dependency and no image file is redistributed with it. Licence: version arXiv:2502.02353v3 of this article is distributed under the Creative Commons Attribution 4.0 International licence, as recorded in the arXiv licence metadata for this exact version and shown on the abstract page https://arxiv.org/abs/2502.02353v3. A full rights notice is delivered at licenses/arxiv-2502.02353-CC-BY-4.0.txt. Characters: every retained excerpt is pure ASCII, so the delivered engine, XeLaTeX with its default Latin Modern text font, reports no missing character.
\begin{Lemma} \label{lem inequ for left-complete}
If the semistable locus \smash{$R_{d_J}^{\tilde{\theta}\text{\rm -sst}} (Q_J)$} is non-empty, then for all left-complete subsets~$\hat{I}$ of the index set~$I$ corresponding to the segment~$J$ the inequality
$\sum_{\hat{I}} \mu_i \leq 0$
holds.
\end{Lemma}

\begin{Theorem} \label{theo balanced semi-stable empty}
For a balanced string $s$ which is not minimal, the semistable locus \smash{$R_{{\bf d}_s}^{\theta_s\text{\rm -sst}}(Q_s)$} is empty.
\end{Theorem}

\begin{proof}
First note that we can switch from the stability $\theta_s$ to the stability $\tilde{\theta}$ given by $\tilde{\theta}(v,j) = -(p_i +j)$ on the level quiver $Q_K$ for a segment $K$ starting with $p_i$, since the open subset of semistable representations does not change when passing from $\theta$ to $x \cdot \theta$ for a positive $x$.

Moreover, we note that for a fixed string $s = (s_1, \dots, s_n) \in \mathbb{Q}^n$, the minimum of $\big\{|s-q| := \sum_{i=1}^n (s_i - q)^2 \big\}$ for $q \in \mathbb{Q}$ is given by $|s-\bar{q}|$ such that $\sum_{i=1}^n (s_i-\bar{q}) =0$. If, additionally, $\sum_{i=1}^n s_i < \sum_{i=1}^n (s_i-q) <0$, then $|s| \geq |s-q|$, since we can consider $|s-q|$ (for a fixed $s$) as a~quadratic function in $q$. Analogously, if $\sum_{i=1}^n s_i > \sum_{i=1}^n (s_i-q) >0$, then $|s| \geq |s-q|$.

Let $\mu = (\mu_1, \dots, \mu_d)$ be a balanced, but not minimal string. By definition, there exists a~(balanced) $\mu'$ with $|\mu'| < |\mu|$ and $\R(\mu') \supset \R(\mu)$. Therefore, we find a segment $J= (\mu_i)_{i \in I}$ in $\mu$ such that the corresponding substring $J' = (\mu')_{i \in I}$ of $\mu'$ satisfies $|J'|= \sum_I \mu_i'^2 < \sum_I \mu_i^2=|J|$.
We can decompose the indexing set $I$ in three disjoint sets given by
\begin{align*}
I^n_1 = \{ i \in I \text{ with } \mu_i > \mu_i' \} ,\qquad
I^p_1 = \{ i \in I \text{ with } \mu_i <\mu_i' \}, \qquad
I_0 = \{ i \in I \text{ with } \mu_i = \mu_i' \} .
\end{align*}
Since $|J'| < |J|$, either $\sum_{I^n_1} \mu_i'^2 < \sum_{I^n_1} \mu_i^2$ or $\sum_{I^p_1} \mu_i'^2 < \sum_{I^p_1} \mu_i^2$. Without loss of generality, we assume the inequality holds for $I^n_1$. The other case follows analogously.

Consider $\mu_i'$ with $i \in I^n_1$. Then there exists an $\epsilon_i >0$ with $\mu_i' = \mu_i - \epsilon_i$. In the level quiver $Q_J$, there exists a vertex $v_i$ corresponding to $\mu_i$. Every arrow $\alpha$ in $Q_J$ with $t(\alpha)= v_i$ corresponds to (possibly several) elements $\mu_l \in J$ such that $\mu_i - \mu_l = 1$ and $\pi_{li} \in \Pi$. Since $\R(\mu') \supset \R(\mu)$, every $\mu_l$ obtained in the above way fulfills $\mu'_l = \mu_l - \epsilon_l$ with $\epsilon_l > \epsilon_i$. This shows that $I^n_1$ is left-complete.

We now assume that the semistable locus in question is not empty.
By Lemma \ref{lem inequ for left-complete}, we know that for all left-complete subset $\hat{I}$ the inequality \smash{$\sum_{\hat{I}} \mu_i \leq 0$} hold.
We know that $\smash{\sum_{I^n_1} \mu_i'^2} = \sum_{I^n_1} (\mu_i - \epsilon_i)^2 < \sum_{I^n_1} \mu_i^2$ and $\sum_{I^n_1} \mu_i \leq 0$, because $I^n_1$ is left-complete. We choose $\tilde{\epsilon}_1$ as the minimum of all $\epsilon_i$ for $i \in I^n_1$. We get
\begin{align*}
\sum_{I^n_1} \mu_i -\epsilon_i < \sum_{I^n_1} \mu_i - \epsilon_i + \tilde{\epsilon}_1 \leq \sum_{I^n_1} \mu_i \leq 0 ,
\end{align*}
leading to \smash{$\sum_{I^n_1} (\mu_i - \epsilon_i)^2 > \sum_{I^n_1} (\mu_i - \epsilon_i + \tilde{\epsilon}_1)^2$}.
Now we define
$
I^n_2 = \{i \in I^n_1 \mid \mu_i -\epsilon_i + \tilde{\epsilon}_1 < \mu_i\} $.
In other words, we exclude all $\mu_i'$ with $\mu_i' + \tilde{\epsilon}_1 = \mu_i$.
Now $I^n_2$ is left-complete by the same argument as for $I^n_1$. We choose $\tilde{\epsilon}_2$ as the minimum of all $\epsilon_i - \tilde{\epsilon}_1$ with $i \in I^n_2$. We get
\begin{align*}
\sum_{I^n_2} \mu_i - \epsilon_i +\tilde{\epsilon}_1 < \sum_{I^n_2} \mu_i - \epsilon_i+ \tilde{\epsilon}_1 + \tilde{\epsilon}_2 \leq \sum_{I^n_2} \mu_i \leq 0
\end{align*}
concluding
\begin{align*}
\sum_{I^n_2} (\mu_i - \epsilon_i + \tilde{\epsilon}_1)^2 > \sum_{I^n_2} (\mu_i - \epsilon_i + \tilde{\epsilon}_1 + \tilde{\epsilon}_2)^2.
\end{align*}
We get a descending chain of subsets
$
I_1^n \supsetneq I_2^n \supsetneq \cdots$,
which has to terminate after finitely many steps. This yields the following estimate:
\begin{align*}
\sum_{I^n_1} (\mu_i -\epsilon_i)^2 > \sum_{I^n_1} (\mu_i - \epsilon_i + \tilde{\epsilon}_1)^2 > \sum_{I^n_1 \setminus I^n_2} \mu_i^2 + \sum_{I^n_2} (\mu_i - \epsilon_i + \tilde{\epsilon}_1 + \tilde{\epsilon}_2)^2 > \dots > \sum_{I^n_1} \mu_i^2 .
\end{align*}
Since we have started with $\sum_{I^n_1} (\mu_i -\epsilon_i)^2 < \sum_{I^n_1} \mu_i^2$, we have thus constructed a contradiction, proving that a balanced, but not minimal string, has an empty semistable locus.
\end{proof}

\end{document}
